\documentclass{report}
\usepackage{amsmath,amssymb}
\setlength{\parindent}{0mm}
\setlength{\parskip}{1em}
\begin{document}
\begin{center}
\rule{6in}{1pt} \
{\large Masahiro Kanazaki \\
{\bf Aircraft's Trajectory Optimization During Descent Using Kriging Model Based Genetic Algorithm}}

Tokyo Metoropolitan University \\ 6-6 Asahigaoka \\ Hino \\ Tokyo \\ Japan
\\
{\tt kana@sd.tmu.ac.jp}\\
Norazila Othman\end{center}

To develop the next generation aircraft, the safety, the efficiency by
reducing fuel consumption and the reducing environmental impact are
critical issues because of higher density landing in the future. A
microburst is a typical hazard to aircraft descent.

Several studies have investigated control strategies to overcome the
hazardous conditions using a successive quadratic programs trajectory
optimization algorithm. They showed the ability of their algorithms to
decide time-series control such as angle of attack, and pitch angle.
However, these studies did not address in detail the major effect of the
aerodynamics associated with the hazardous conditions during takeoff or
descent.

In this study, the time-series trajectory optimization problem for a
conventional civil aircraft is considered. The trajectory is estimated by
the 3-degree-of-freedom (3DoF) equations of motion (EoM); the
aerodynamics are estimated through the trajectory evaluation by USAF
stability and control DATCOM. Global optimization is carried out by means
of Kriging model based multi-objective evolutionary algorithm (MOEA).
Objective functions are to minimize the cost function, which indicates
the efficiency trajectory profile during the descent of an aircraft and
the minimization of the maximum acceleration, which decides the impact
load to passengers and payloads. The trajectory is evaluated time wise
and two cases are compared; one is the trajectory when no microburst
occurs and the other is the trajectory with a microburst. Analysis of
variance (ANOVA) is used to investigate effects of the input variables.


\end{document}
